\subsection{超度量绝对值}

由 K 到 \(R_{+}\) 的映射 f 若满足条件 ( \(VA_{I}\) )、( \(VA_{II}\) ) 与 ( \(U_{1}\) ) （这显然蕴含 f 是绝对值），则称为超度量绝对值。

命题 3. 设 \(f\) 是由 \(\mathbf{K}\) 到 \(\mathbf{R}_+\) 的映射。下列性质等价：

\begin{enumerate}
\def\labelenumi{(\alph{enumi})}
\item
  \(f\) 是超度量绝对值。
\item
  存在一个赋值 \(v\) （在 \(\mathbf{K}\) 上取值于 \(\mathbf{R}\) ）以及实数 \(a\) ，使得 \(0 < a < 1\) 且 \(f = a^v\) 。
\item
  \(f\) 属于 \(\mathcal{V}(\mathbf{K})\) ，且 \(f(n,1) \leqslant 1\) 对每个整数 \(n > 0\) 成立。
\item
  对一切 \(s > 0, f^s\) 是绝对值。
\end{enumerate}

对每个实数 \(c\) （满足 \(0 < c < 1\) ），映射 \(t \mapsto c^t\) 是有序群 \(\mathbf{R}\) （带与通常序相反的序）到有序群 \(\mathbf{R}_+^*\) 上的同构；这表明 (a) 与 (b) 等价。显然 (a) 蕴含 (c)；(c) 蕴含 (d)，因为由 (c) 推出

\[
(f (n. 1)) ^ {s} \leqslant 1 \leqslant n
\]

对每个整数 \(n > 0\) 成立，而第 1 节命题 2 表明 \(f^s\) 是绝对值。最后 (d) 蕴含 (a)：因为若 \(f^s\) 是绝对值，它就满足 \((\mathbf{U}_2)\) ，因而 \(f\) 满足 \((\mathbf{U}_{2^{1 / s}})\) 对一切 \(s > 0\) 成立，因此它也满足 \((\mathbf{U}_1)\) ，令 \(s\) 趋于 \(+\infty\) 。

推论. 若 K 是特征 p > 0 的（未必交换的）域，则 \(\mathcal{V}(\mathbf{K})\) 上的每个函数都是超度量绝对值。

每个非零元素 \(z = n.1\) （ \(n\) 为整数 \(>0\) ）都属于素子域 \(\mathbf{F}_p\) （\(K\) 的），因而满足关系 \(z^{p-1} = 1\) ，这蕴含 \(f(z) = 1\) ，于是可以应用命题 3 (c)。

给定实数 \(c\) （满足 \(0 < c < 1\) ），公式

\[
f (x) = c ^ {v (x)}, \quad v (x) = \log_ {c} f (x)
\]

就在 K 上的超度量绝对值与 K 上取实值的赋值之间建立了一一对应。非真赋值对应于非真绝对值（《一般拓扑学》第九章 § 3 第 2 节）。设 \(v_{1}, v_{2}\) 是 K 上两个取实值的赋值，\(f_{1}, f_{2}\) 是相应的绝对值；\(v_{1}\) 与 \(v_{2}\) 等价，当且仅当 \(f_{1}\) 与 \(f_{2}\) 等价：因为说 \(v_{1}\) 与 \(v_{2}\) 等价，相当于说关系 \(v_{1}(x) \geqslant 0\) 与 \(v_{2}(x) \geqslant 0\) 等价，或者再说关系 \(f_{1}(x) \leqslant 1\) 与 \(f_{2}(x) \leqslant 1\) 等价；因此只需应用《一般拓扑学》第九章 §3 第 2 节命题 5。此外（同前），由 \(f_{1}\) 与 \(f_{2}\) 在 K 上定义的拓扑相同，当且仅当 \(f_{1}\) 与 \(f_{2}\) 等价。

